% GENERATED FILE. Edit canonical lessons, then run npm run generate:books.
% canonical-source-hash: fnv1a64:f6896faa6a9859bc
% canonical-lessons: RU-C40-spat, RU-C40-sidet, RU-C40-stoyat, RU-C40-bezhat, RU-C40-igrat

\chapter{Sleep, Sit, Stand, Run, Play}
\label{ch:ru-sleep-sit-stand-run-play}
\hlchaptermodality{40}

\begin{chapteropening}
\textbf{By the end of this chapter:} \emph{I can say to sleep, to sit, to stand, to run and to play.}

Complete the last lesson of chapter 40: I can say to sleep, to sit, to stand, to run and to play.
\end{chapteropening}

\section[\texorpdfstring{\ru{спать} (spat)}{спать (spat)}]{\ru{спать} (spat) — to sleep}

\label{lesson:RU-C40-spat}



\begin{quote}
Before the new one: say the Russian for to sell, then the Russian for to swim.
\end{quote}

\subsection*{You'll want to know: \ru{спать}}
\textbf{\ru{спать}} (\emph{spat}) — \textquotedblleft{}to sleep\textquotedblright{}.

\begin{cousinweb}
\textbf{\ru{спать}} $\leftarrow$ Proto-Slavic \emph{*sŭpati} — inherited from PIE \emph{*swep-}, \textquotedblleft{}to sleep\textquotedblright{}. Vasmer lists Sanskrit \emph{svápiti}, \textquotedblleft{}sleeps\textquotedblright{}, Old English \emph{swefan}, \textquotedblleft{}to sleep\textquotedblright{}, and Latin \emph{sōpiō}, \textquotedblleft{}I lull to sleep\textquotedblright{}.
\end{cousinweb}

\begin{grammarlens}[title={What you've built}]
One more word for this chapter.
\end{grammarlens}

\subsection*{Your turn}
\begin{itemize}
\raggedright
  \item \emph{Say it:} \emph{spat}
  \item \emph{Say it:} \emph{spat}, once more
  \item \emph{Say it:} say \emph{spat}
  \item \emph{Recall:} say the Russian for to sell, then the Russian for to swim, then say \emph{spat} again
\end{itemize}

\begin{tcolorbox}[breakable,colback=teal!4,colframe=teal!35!black,title={Before you move on}]
What does \ru{спать} mean? (\textquotedblleft{}To sleep\textquotedblright{}.) Say it once more.
\end{tcolorbox}
\section[\texorpdfstring{\ru{сидеть} (sidét)}{сидеть (sidét)}]{\ru{сидеть} (sidét) — to sit}

\label{lesson:RU-C40-sidet}



\begin{quote}
Before the new one: say the Russian for to swim, then the Russian for to sleep.
\end{quote}

\subsection*{You'll want to know: \ru{сидеть}}
\textbf{\ru{сидеть}} (\emph{sidét}) — \textquotedblleft{}to sit\textquotedblright{}.

\begin{cousinweb}
\textbf{\ru{сидеть}} $\leftarrow$ Proto-Slavic \emph{*sěděti} — inherited from PIE \emph{*sed-}, \textquotedblleft{}to sit\textquotedblright{}. Vasmer lists Lithuanian \emph{sėdėti}, Latin \emph{sedeō} and Gothic \emph{sitan}: English \emph{sit} is a cousin.
\end{cousinweb}

\begin{grammarlens}[title={What you've built}]
One more word for this chapter.
\end{grammarlens}

\subsection*{Your turn}
\begin{itemize}
\raggedright
  \item \emph{Say it:} \emph{sidét}
  \item \emph{Say it:} \emph{sidét}, once more
  \item \emph{Say it:} say \emph{sidét}
  \item \emph{Recall:} say the Russian for to swim, then the Russian for to sleep, then say \emph{sidét} again
\end{itemize}

\begin{tcolorbox}[breakable,colback=teal!4,colframe=teal!35!black,title={Before you move on}]
What does \ru{сидеть} mean? (\textquotedblleft{}To sit\textquotedblright{}.) Say it once more.
\end{tcolorbox}
\section[\texorpdfstring{\ru{стоять} (stoyát)}{стоять (stoyát)}]{\ru{стоять} (stoyát) — to stand}

\label{lesson:RU-C40-stoyat}



\begin{quote}
Before the new one: say the Russian for to sleep, then the Russian for to sit.
\end{quote}

\subsection*{You'll want to know: \ru{стоять}}
\textbf{\ru{стоять}} (\emph{stoyát}) — \textquotedblleft{}to stand\textquotedblright{}.

\begin{cousinweb}
\textbf{\ru{стоять}} $\leftarrow$ Proto-Slavic \emph{*stojati} — inherited from PIE \emph{*steh\textsubscript{2}-}, \textquotedblleft{}to stand\textquotedblright{}. Vasmer compares Oscan \emph{staít}, \textquotedblleft{}stands\textquotedblright{}, Latin \emph{stāre} and Greek \emph{hístēmi}, \textquotedblleft{}I set up\textquotedblright{}: English \emph{stand} is a cousin too.
\end{cousinweb}

\begin{grammarlens}[title={What you've built}]
One more word for this chapter.
\end{grammarlens}

\subsection*{Your turn}
\begin{itemize}
\raggedright
  \item \emph{Say it:} \emph{stoyát}
  \item \emph{Say it:} \emph{stoyát}, once more
  \item \emph{Say it:} say \emph{stoyát}
  \item \emph{Recall:} say the Russian for to sleep, then the Russian for to sit, then say \emph{stoyát} again
\end{itemize}

\begin{tcolorbox}[breakable,colback=teal!4,colframe=teal!35!black,title={Before you move on}]
What does \ru{стоять} mean? (\textquotedblleft{}To stand\textquotedblright{}.) Say it once more.
\end{tcolorbox}
\section[\texorpdfstring{\ru{бежать} (bezhát)}{бежать (bezhát)}]{\ru{бежать} (bezhát) — to run}

\label{lesson:RU-C40-bezhat}



\begin{quote}
Before the new one: say the Russian for to sit, then the Russian for to stand.
\end{quote}

\subsection*{You'll want to know: \ru{бежать}}
\textbf{\ru{бежать}} (\emph{bezhát}) — \textquotedblleft{}to run\textquotedblright{}.

\begin{cousinweb}
\textbf{\ru{бежать}} $\leftarrow$ Proto-Slavic \emph{*běžati} — inherited, a cousin of Lithuanian \emph{bėgti}, \textquotedblleft{}to run\textquotedblright{}, and of Greek \emph{phébomai}, \textquotedblleft{}I flee\textquotedblright{}, and \emph{phóbos}, \textquotedblleft{}flight, fear\textquotedblright{} (Vasmer), the source of English \emph{phobia}.
\end{cousinweb}

\begin{grammarlens}[title={What you've built}]
One more word for this chapter.
\end{grammarlens}

\subsection*{Your turn}
\begin{itemize}
\raggedright
  \item \emph{Say it:} \emph{bezhát}
  \item \emph{Say it:} \emph{bezhát}, once more
  \item \emph{Say it:} say \emph{bezhát}
  \item \emph{Recall:} say the Russian for to sit, then the Russian for to stand, then say \emph{bezhát} again
\end{itemize}

\begin{tcolorbox}[breakable,colback=teal!4,colframe=teal!35!black,title={Before you move on}]
What does \ru{бежать} mean? (\textquotedblleft{}To run\textquotedblright{}.) Say it once more.
\end{tcolorbox}
\section[\texorpdfstring{\ru{играть} (igrát)}{играть (igrát)}]{\ru{играть} (igrát) — to play}

\label{lesson:RU-C40-igrat}



\begin{quote}
Before the new one: say the Russian for to stand, then the Russian for to run.
\end{quote}

\subsection*{You'll want to know: \ru{играть}}
\textbf{\ru{играть}} (\emph{igrát}) — \textquotedblleft{}to play\textquotedblright{}.

\begin{cousinweb}
\textbf{\ru{играть}} goes with \emph{\ru{игра}}, \textquotedblleft{}a game\textquotedblright{}, $\leftarrow$ Proto-Slavic \emph{*jĭgra}, \emph{*jĭgrati} — inherited. Vasmer relates it to Lithuanian \emph{áikštytis}, \textquotedblleft{}to be playful\textquotedblright{}, and Sanskrit \emph{ejati}, \textquotedblleft{}stirs, moves\textquotedblright{}.
\end{cousinweb}

\begin{grammarlens}[title={What you've built}]
One more word for this chapter.
\end{grammarlens}

\subsection*{Your turn}
\begin{itemize}
\raggedright
  \item \emph{Say it:} \emph{igrát}
  \item \emph{Say it:} \emph{igrát}, once more
  \item \emph{Say it:} say \emph{igrát}
  \item \emph{Recall:} say the Russian for to stand, then the Russian for to run, then say \emph{igrát} again
\end{itemize}

\begin{tcolorbox}[breakable,colback=teal!4,colframe=teal!35!black,title={Before you move on}]
What does \ru{играть} mean? (\textquotedblleft{}To play\textquotedblright{}.) Say it once more.
\end{tcolorbox}
